\chapter{Security and untrusted input}
\label{ch:security}

\begin{aside}
A TUI app processes input from many sources: the user's
keyboard, the user's clipboard, files the user opens, network
data the app fetches, output from subprocesses. Any of these
can be hostile. This chapter is a short tour of the patterns
\maya{} encourages for keeping the app safe.
\end{aside}

\section{What \emph{can} go wrong}

A handful of failure modes are specific to terminal apps:

\begin{itemize}[leftmargin=*]
  \item \textbf{ANSI injection.} Content that contains escape
    codes can change the terminal's state in ways your code
    did not authorise. \code{\\e[2J} clears the screen;
    \code{\\e]52;c;\ldots\\e\\textbackslash} writes the
    clipboard.
  \item \textbf{Format string vulnerabilities.} If you pass
    untrusted text to \code{printf}-like functions as the
    format string, you have an information leak.
  \item \textbf{Buffer overflows in C interop.} If you call C
    code with the wrong size, you can corrupt memory.
  \item \textbf{Subprocess injection.} If you build a command
    line by concatenating user input, you create shell-
    injection paths.
  \item \textbf{Path traversal.} If you open a file by
    user-supplied path, the user can read anywhere your
    process has rights.
\end{itemize}

Each has a specific mitigation.

\section{ANSI injection}

The classic. A log line containing
\code{\\e[2J\\e[H} clears the screen and homes the cursor.
A pasted blob with \code{\\e]52;c;<base64>\\e\\textbackslash}
writes the clipboard.

\subsection*{Mitigation: strip on display}

\maya{} provides \code{sanitize\_text}:

\begin{lstlisting}
std::string sanitize_text(std::string_view input) {
    std::string out;
    out.reserve(input.size());
    for (size_t i = 0; i < input.size(); ++i) {
        unsigned char c = input[i];
        if (c == 0x1B) { i++; while (i < input.size() && !is_terminator(input[i])) i++; continue; }
        if (c < 0x20 && c != '\\t' && c != '\\n') continue;
        out += c;
    }
    return out;
}
\end{lstlisting}

Strip every escape sequence; strip every control character
except tab and newline. The result is safe to display.

\maya{} runs \code{sanitize\_text} on:

\begin{itemize}[leftmargin=*]
  \item Pasted text (from bracketed paste).
  \item File contents you display (when wrapped in
    \code{text\_safe}).
  \item Network-fetched strings via \code{Cmd::task}.
\end{itemize}

\subsection*{Mitigation: lock down the terminal}

Some terminals support a ``deny mode'' for clipboard writes
(tmux \code{set-clipboard off}). Document this for users who
display untrusted content frequently.

\section{Format string safety}

C++26's \code{std::format} is type-safe: the format string must
match the argument count and types, checked at compile time:

\begin{lstlisting}
std::format(\"{}\", value);                     // OK
std::format(input, value);                     // ERROR
\end{lstlisting}

The second form does not compile because \code{input} is not a
literal. Use \code{std::vformat} only when you genuinely need
a runtime format string; even then, validate it first.

\maya{} provides a wrapper that rejects suspicious format
strings (containing \code{\%n}, etc.). For ordinary use,
\code{std::format} is the safe path.

\section{Subprocess injection}

\code{std::system} and \code{popen} are unsafe with user input:

\begin{lstlisting}
std::system(\"grep \" + user_pattern + \" file\");  // injection!
\end{lstlisting}

If \code{user\_pattern} is \code{;rm -rf /}, the result is a
disaster.

\subsection*{Mitigation: argv-based exec}

\maya{} ships an exec helper that takes a vector, not a string:

\begin{lstlisting}
auto result = exec(\"grep\", {user_pattern, \"file\"});
\end{lstlisting}

The shell is never involved. Each element of the vector is one
argv slot. No quoting, no escaping, no injection.

On POSIX this maps to \code{execvp}; on Win32 to
\code{CreateProcessW} with the argv re-encoded.

\section{Path traversal}

\code{file::read(\"~/data/\" + filename)} where \code{filename}
is \code{../../etc/passwd} is a problem.

\subsection*{Mitigation: canonicalise and verify}

\begin{lstlisting}
std::filesystem::path requested = base / filename;
std::filesystem::path canonical = std::filesystem::weakly_canonical(requested);
if (!canonical.string().starts_with(base.string())) {
    return std::unexpect, IoError{\"path escape\"};
}
\end{lstlisting}

\code{weakly\_canonical} resolves \code{..} segments. After
resolution, check that the result lives under the intended
base directory.

\section{Untrusted serialised data}

A JSON or TOML file the user supplies could be malicious:

\begin{itemize}[leftmargin=*]
  \item Deeply nested objects can stack-overflow a recursive
    parser.
  \item Large arrays can exhaust memory.
  \item Specific patterns can trigger expensive regex matches
    (\emph{ReDoS}).
\end{itemize}

Use a strict parser with explicit limits:

\begin{lstlisting}
JsonOptions opt;
opt.max_depth = 64;
opt.max_array_size = 100'000;
opt.max_string_size = 1'000'000;

auto r = parse_json(input, opt);
\end{lstlisting}

The parser rejects content that exceeds limits with a clear
error.

\section{Memory safety}

C++ does not give you memory safety by default. Patterns that
help:

\begin{itemize}[leftmargin=*]
  \item Use \code{std::string\_view}, \code{std::span},
    \code{std::vector}, \code{std::array} instead of raw
    pointers and lengths.
  \item Enable sanitizers (\code{-fsanitize=address,undefined})
    in debug builds.
  \item Use \code{static\_assert} liberally to catch size
    mistakes at build.
  \item Bounds-check in places that touch external data
    (parsers, network, file readers).
\end{itemize}

\maya{} is internally written in this style. External callers
should follow suit.

\section{Cryptography}

\maya{} does not ship crypto. If your app needs it:

\begin{itemize}[leftmargin=*]
  \item Use \code{libsodium} for general-purpose crypto.
  \item Use the OS's TLS library for HTTPS (\code{libcurl}
    + system trust store).
  \item Do not roll your own. Ever.
\end{itemize}

\section{Privilege management}

If your app runs as root (unusual for a TUI), drop privileges
as early as possible:

\begin{itemize}[leftmargin=*]
  \item \code{setuid}/\code{setgid} to a non-privileged user.
  \item \code{chroot} into a restricted directory if you only
    need a subtree.
  \item Use \code{seccomp} or platform sandboxing to restrict
    system calls.
\end{itemize}

For ordinary apps, the user already runs you with their own
privileges. The risk is in what you do with those privileges.

\section{Logging secrets}

Already covered (Chapter~\ref{ch:observability}). Redact:

\begin{itemize}[leftmargin=*]
  \item API tokens.
  \item Passwords (never log these in any form).
  \item Personal data (emails, addresses).
  \item File contents in error messages (a path is fine; the
    body is not).
\end{itemize}

\section{Supply chain}

Every dependency is an attack surface. Mitigations:

\begin{itemize}[leftmargin=*]
  \item Pin versions in the build system; do not float to
    latest.
  \item Vendor critical dependencies (commit them to your
    repository).
  \item Audit periodically (CVE databases, \code{cargo audit}
    equivalents for C++).
  \item Limit the dependency surface --- prefer the standard
    library where possible.
\end{itemize}

\maya{} itself has a small dependency footprint: the standard
library, optional toml++ for config, optional xxHash for
fingerprints. Each is auditable.

\section{Pitfalls}

\begin{warning}
\textbf{Displaying raw subprocess output.} If \code{git diff}
returns content with ANSI codes (which it does for colour),
those codes will be interpreted by the terminal. Sanitize
first.
\end{warning}

\begin{warning}
\textbf{Writing user-controlled text to the clipboard.} If a
user pastes a malicious blob and you write it to the clipboard
unmodified, the next paste somewhere else can carry the
payload. Validate before writing.
\end{warning}

\begin{warning}
\textbf{Trusting file paths.} A symlink can point anywhere.
Resolve fully before opening.
\end{warning}

\begin{warning}
\textbf{Verbose error messages with paths.} A bug report
including the absolute path of a user's home directory leaks
information. Print relative paths or anonymised paths.
\end{warning}

\section*{Workshop}

\begin{enumerate}[leftmargin=*]
  \item Add \code{sanitize\_text} to a log viewer. Test with a
    file containing ANSI escapes; verify they are stripped.
  \item Replace a \code{std::system} call with the argv-based
    \code{exec} helper. Confirm the same behaviour without
    shell involvement.
  \item Audit your dependency list. For each, note the last
    update and any known CVEs.
\end{enumerate}

\begin{recap}
Untrusted input is everywhere: keyboard, clipboard, files,
network, subprocesses. Each has a mitigation: sanitize ANSI
escapes, use type-safe formatting, argv-based exec, canonical
path checks, strict parsers with limits, redacted logging. The
mitigations are small individually; together they keep the app
out of the news.
\end{recap}
